\documentclass[11pt]{article}
\usepackage{amsmath,amssymb,booktabs,geometry,graphicx,hyperref,url}
\usepackage{subcaption,placeins}
\hypersetup{hidelinks}
\geometry{margin=1in}
\setlength{\parskip}{0.45em}
\setlength{\parindent}{0pt}
\graphicspath{{./}{../}}
\title{Stage 8: Positive Continuation-Premium Price Surrogate}
\author{Free-Boundary-Aware Risk Surfaces for American Options}
\date{\today}
\begin{document}
\maketitle

\section{Stage 8: Positive Continuation-Premium Price Surrogate}

\subsection{Purpose}
This stage evaluates obstacle-consistent price and continuation-premium
surrogates. Instead of predicting price directly, the model predicts the
nonnegative value of waiting rather than exercising and reconstructs price as
payoff plus continuation premium. This removes the specific failure observed in
Stage 7, but it remains price/premium evidence only: it is not boundary, Delta,
Gamma, or PSOR-replacement evidence.

\subsection{Method}
The clean Stage 6 dataset is used with its regime-level train, validation, and test splits. Two positive-premium models are audited: a linear sanity baseline trained on $\log(1+\mathrm{premium})$, and a neural positive-premium model using a softplus output. Both models reconstruct price by adding predicted premium to the payoff.

\subsection{Mathematical Formulation}
Let $\Phi(S)$ denote the payoff and $U(S,\tau)$ the CN/PSOR American option value. The continuation premium is
\[
C(S,\tau)=U(S,\tau)-\Phi(S)\geq 0.
\]
The surrogate predicts
\[
\widehat{C}_{\theta}(x)=\operatorname{softplus}(f_{\theta}(x))\geq 0,
\]
and reconstructs
\[
\widehat{U}_{\theta}(S,\tau)=\Phi(S)+\widehat{C}_{\theta}(x).
\]
Therefore $\widehat{U}_{\theta}\geq \Phi$ by construction for the neural positive-premium model.

\subsection{Implementation Details}
The training workflow is implemented in the Stage 8 positive-premium experiment. The delicate-region error profiles are implemented in the companion boundary-distance profile script. The model factory and reconstruction rule are in \texttt{src/models/premium\_surrogate.py}.

The run used capped deterministic samples: 250,000 train samples, 100,000 validation samples, and 250,000 test samples. The selected neural artifact is written under \texttt{results/models}. The linear model is retained only as a sanity baseline.

\subsection{Validation and Diagnostics}
\begin{table}[h]
\centering
\begin{tabular}{llrrrr}
\toprule
Model & Split & Price RMSE & Price MAE & Premium RMSE & Max obstacle\\
\midrule
Linear premium & Validation & 0.0188 & 0.0104 & 0.0188 & 0.0000\\
Linear premium & Test & 0.0201 & 0.0106 & 0.0201 & 0.0000\\
MLP premium & Validation & 0.0200 & 0.0098 & 0.0200 & 0.0000\\
MLP premium & Test & 0.0218 & 0.0108 & 0.0218 & 0.0000\\
\bottomrule
\end{tabular}
\caption{Validation and test price/premium errors for the linear sanity model and softplus MLP. Both reconstruct \(\widehat U=\Phi+\widehat C\), so the payoff obstacle is satisfied on evaluated samples by construction.}
\end{table}

\begin{table}[h]
\centering
\begin{tabular}{llrr}
\toprule
Model & Split & Boundary-band RMSE & Strict-mask RMSE\\
\midrule
Linear premium & Validation & 0.0098 & 0.0175\\
Linear premium & Test & 0.0094 & 0.0190\\
MLP premium & Validation & 0.0055 & 0.0187\\
MLP premium & Test & 0.0062 & 0.0207\\
\bottomrule
\end{tabular}
\caption{Price RMSE in the free-boundary band and strict-mask region. The softplus MLP improves boundary-band error but does not outperform the linear sanity baseline in aggregate RMSE.}
\end{table}

\subsection{Results}
Both positive-premium models have zero maximum obstacle violation on evaluated
samples. The linear sanity baseline has the lower aggregate validation RMSE,
0.0188. The selected neural MLP has validation RMSE 0.0200 and test RMSE
0.0218; it is retained as the neural artifact because it implements the planned
softplus parameterization and has lower boundary-band error. This selection is
not evidence that the MLP dominates the linear model.

\begin{figure}[h]
\centering
\includegraphics[width=0.68\linewidth]{results/figures/stage08_positive_premium_surrogate/validation_price_rmse.png}
\caption{Validation price RMSE for the linear and neural positive-premium models. The linear sanity baseline has the lower aggregate RMSE in this run.}
\end{figure}

The maximum test-set obstacle violation is zero for both positive-premium
models. Because an all-zero bar chart would add no visual information, this
result is reported in the metric table and explained by the reconstruction
\(\widehat U=\Phi+\widehat C\).

Boundary-distance diagnostics for the selected neural model show test RMSE 0.0061 near the boundary, 0.0220 away from the boundary, 0.0844 near the payoff kink, and 0.0208 under the strict mask. The high near-kink error reinforces the earlier conclusion that nonsmooth regions require separate reporting.

% Auto-generated by experiments/19_report_visualization_enrichment.py
\section*{Visual Evidence Gallery}
\begin{figure}[!htbp]
\centering
\begin{subfigure}[t]{0.48\linewidth}
\centering
\includegraphics[width=\linewidth]{results/figures/report_enrichment/stage08_positive_premium_surrogate_positive_premium_architecture.pdf}
\caption{Softplus premium output enforces nonnegative continuation premium before price reconstruction.}
\end{subfigure}
\hfill
\begin{subfigure}[t]{0.48\linewidth}
\centering
\includegraphics[width=\linewidth]{results/figures/report_enrichment/stage08_positive_premium_surrogate_price_reconstruction.pdf}
\caption{Price is reconstructed as payoff plus nonnegative premium and remains a price only model.}
\end{subfigure}
\end{figure}

\begin{figure}[!htbp]
\centering
\begin{subfigure}[t]{0.48\linewidth}
\centering
\includegraphics[width=\linewidth]{results/figures/report_enrichment/stage08_positive_premium_surrogate_premium_metrics.pdf}
\caption{Positive-premium models are compared on price and premium errors under regime-level splits; the linear sanity baseline retains lower aggregate RMSE in this run.}
\end{subfigure}
\hfill
\begin{subfigure}[t]{0.48\linewidth}
\centering
\includegraphics[width=\linewidth]{results/figures/report_enrichment/stage08_positive_premium_surrogate_boundary_distance_profile.pdf}
\caption{Boundary distance profiles show where the price surrogate is weak without treating it as a boundary model.}
\end{subfigure}
\end{figure}

\FloatBarrier


\subsection{Failure Modes and Limitations}
This stage does not prove that the neural MLP dominates the linear sanity model; the linear baseline has lower aggregate validation and test RMSE in this run. The positive-premium construction guarantees obstacle consistency, but it does not guarantee monotonicity, convexity, boundary accuracy, Delta accuracy, or Gamma accuracy. The near-kink error remains large, and the selected model is scoped to the clean Stage 6 regimes.

\subsection{Conclusion and Next Step}
The Stage 8 gate passes for obstacle-consistent price/premium modeling within the tested regime envelope. The next stage should train a direct boundary diagnostic head and evaluate it separately from this price surrogate.

\end{document}
